\chapter{Elaboração de questões em Matemática: referenciais didáticos}
\label{cap:didatica}

% =====================================================================
% PAPEL DESTE CAPÍTULO (registro didático):
%   Responde "o que é uma questão boa e como sei disso".
%   Cada referencial DERIVA implicações; a lista formatada dos SEIS
%   critérios de qualidade existe APENAS na Seção 3.4 (casa canônica).
%   O par "TBR na entrada, SOLO na avaliação" é dito UMA vez, em 3.2.3.
% =====================================================================

A arquitetura multiagente descrita no capítulo anterior sustenta-se sobre dois conceitos que ainda
carecem de fundamentação explícita: a especificação de entrada que o professor fornece ao sistema e a
rubrica utilizada pelo agente crítico didático para avaliar as questões geradas. Ambos precisam estar
ancorados em referenciais teóricos da educação matemática reconhecidos pela comunidade acadêmica e
coerentes com os objetivos do Ensino Médio brasileiro. Este capítulo apresenta, em quatro seções
articuladas, os pilares pedagógicos do sistema --- a perspectiva de Pólya; a Teoria das Situações
Didáticas de Brousseau; as taxonomias cognitivas de Bloom e SOLO; e as habilidades de Matemática da
BNCC do Ensino Médio --- e encerra com a consolidação desses referenciais em critérios de qualidade
operacionalizáveis (Seção~\ref{sec:criterios}).

\section{Resolução de problemas e elaboração de questões}
\label{sec:resolucao}

A resolução de problemas ocupa posição central no ensino de Matemática há décadas. No Brasil, essa
centralidade está expressa tanto nos Parâmetros Curriculares Nacionais \citep{brasil1998pcn} quanto na
Base Nacional Comum Curricular \citep{brasil2018bncc}, que reafirmou o desenvolvimento do raciocínio
lógico e da capacidade de resolução de situações-problema como competências essenciais da área.
Antes de apresentar a perspectiva de Pólya, é necessário estabelecer uma distinção que atravessa toda
a tradição do campo: a diferença entre \textbf{exercício} e \textbf{problema}.

Um exercício é uma tarefa cuja via de resolução é imediatamente reconhecível pelo estudante: basta
aplicar um algoritmo já aprendido, sem necessidade de decisão estratégica. Um problema, por sua vez, é
uma situação em que o resolvedor não dispõe, de imediato, de um caminho garantido para a solução;
existe uma lacuna entre o que se sabe e o que se quer determinar, e superá-la exige planejamento e
raciocínio. Para \citet{dante2003didatica}, o exercício é essencialmente uma atividade de
\emph{aplicação} de habilidade conhecida, ao passo que o problema necessariamente envolve invenção ou
criação significativa por parte de quem o resolve. Essa distinção não é absoluta: uma mesma tarefa pode
ser exercício para quem domina o procedimento e problema genuíno para quem ainda não o construiu --- o
que define o caráter de um problema não é sua dificuldade intrínseca, mas o grau de mobilização
cognitiva que ele exige.

\subsection{A perspectiva de Pólya}
\label{sub:polya}

Em 1945, o matemático húngaro George Pólya publicou \emph{How to Solve It}
\citep{polya1945howtosolveit}, obra que se tornaria referência fundamental para a educação matemática
e que no Brasil circula principalmente pela tradução \emph{A Arte de Resolver Problemas}
\citep{polya1995arte}. A proposta central é que o processo de resolução pode ser descrito, ensinado e
aprendido por meio de uma estrutura heurística de quatro etapas.

A primeira etapa é \textbf{compreender o problema}: identificar com clareza o que é dado, o que é
pedido e quais condições conectam os dados à incógnita. A segunda é \textbf{elaborar um plano} ---
estabelecer a conexão entre os dados e a solução, reconhecendo um problema análogo, introduzindo
notação auxiliar ou decompondo o problema em partes; é nessa etapa que reside o núcleo criativo da
resolução. A terceira é \textbf{executar o plano}, colocando a estratégia em prática e verificando cada
passo. A quarta é \textbf{examinar a solução obtida}: ao revisitar a resolução, o estudante verifica se
o resultado é plausível e se a abordagem pode ser generalizada --- razão pela qual questões bem
elaboradas devem admitir uma solução verificável por via independente.

\paragraph{Implicações para a elaboração de questões.} A heurística de Pólya, embora concebida como
guia para o resolvedor, implica exigências precisas para a elaboração de questões formativas. Para que
o estudante possa \emph{compreender o problema}, o enunciado precisa apresentar os dados, o que é
pedido e as condições sem ambiguidades (\emph{clareza do enunciado}). Para que precise \emph{elaborar um
plano}, o enunciado não pode sinalizar tão explicitamente o procedimento que dispense a decisão
estratégica (\emph{não-trivialidade do caminho de solução}). E a etapa de \emph{exame da solução}
pressupõe que a resposta possa ser examinada e confirmada por via independente (\emph{determinação e
verificabilidade da solução}). Essas três exigências são consolidadas, junto com as demais, na
Seção~\ref{sec:criterios}; elas também fundamentam, de forma simétrica, a necessidade do Agente
Verificador Simbólico, que realiza de forma determinística a verificação passo a passo para as classes
de questões tratáveis simbolicamente.

\subsection{Teoria das Situações Didáticas}
\label{sub:tsd}

A Teoria das Situações Didáticas (TSD), desenvolvida por Guy Brousseau a partir da década de 1970
\citep{brousseau1997tsd}, oferece um quadro complementar ao de Pólya. Enquanto a heurística de Pólya
descreve o processo cognitivo do resolvedor individual, a TSD descreve a estrutura da situação que o
professor projeta: o ambiente de aprendizagem e as condições que tornam possível a construção do saber
pelo próprio estudante.

O conceito fundante é o de \textbf{situação didática}, o conjunto de relações estabelecidas entre o
estudante, um ambiente e um sistema educativo, com a intenção de que o estudante se aproprie de um
saber. Dentro dessa estrutura, o conceito de \textbf{\emph{milieu}} é central: tudo aquilo que age
sobre o estudante ou que é acionado por ele --- o problema proposto, os dados fornecidos, as respostas
já obtidas. O \emph{milieu} funciona como um sistema antagonista que oferece resistência e
retroalimentação às ações do estudante, permitindo que ele verifique, por seus próprios meios, se suas
decisões são adequadas. É precisamente essa resistência que torna a aprendizagem possível.

O terceiro conceito central é o \textbf{contrato didático}: o conjunto de expectativas mútuas, em sua
maior parte implícitas, que regulam a relação pedagógica. Brousseau identifica nele efeitos perversos
de relevância direta para a elaboração de questões, sendo o mais conhecido o \textbf{efeito Topaze}:
diante de um estudante em dificuldade, o professor tende a oferecer pistas cada vez mais explícitas até
que a resposta se torne evidente, transferindo para si a responsabilidade intelectual que deveria ser
do estudante. Transposto para a elaboração de questões, esse efeito aparece em enunciados
excessivamente detalhados que conduzem o estudante à solução por um caminho previamente pavimentado.

Quanto à tipologia, Brousseau descreve quatro tipos de situação pelos quais uma sequência didática bem
projetada deve passar: \emph{ação} (o estudante age mobilizando estratégias ainda implícitas),
\emph{formulação} (precisa comunicar suas estratégias, explicitando-as em linguagem), \emph{validação}
(deve convencer um interlocutor produzindo argumentos próximos da demonstração) e
\emph{institucionalização} (o professor formaliza o saber construído, conectando-o ao conhecimento
estabelecido). É necessário reconhecer que a TSD foi concebida para descrever sequências didáticas
completas, e não itens avaliativos isolados; aplicar sua tipologia diretamente a uma questão individual
seria um forçamento teórico. O que a teoria oferece a este trabalho é um vocabulário preciso para
articular o que torna uma questão pedagogicamente rica: uma boa questão constrói um \emph{milieu}
informativo e resistente --- fornece os dados necessários para agir, mas não antecipa a estratégia ---
e evita o efeito Topaze. Essas implicações se traduzem em dois critérios --- a adequação do
\emph{milieu} e a ausência de condução excessiva ---, consolidados na Seção~\ref{sec:criterios}.

\section{Taxonomias cognitivas}
\label{sec:taxonomias}

Os referenciais da seção anterior oferecem critérios qualitativos para avaliar uma questão, mas não
fornecem uma escala que permita \emph{graduar} sua complexidade de forma sistemática. Para isso, o
campo desenvolveu as \textbf{taxonomias cognitivas}: classificações hierárquicas dos processos mentais
exigidos por diferentes tipos de tarefa.

\subsection{Taxonomia de Bloom revisada}
\label{sub:bloom}

A taxonomia original de objetivos educacionais, proposta por Benjamin Bloom e colaboradores em 1956,
foi revisada em 2001 sob a liderança de Lorin Anderson \citep{anderson2001taxonomy}. É essa versão
revisada, conhecida como Taxonomia de Bloom Revisada (TBR), que é adotada nesta dissertação. Sua
principal inovação é a introdução de um \emph{framework} bidimensional: a \textbf{dimensão do
conhecimento} (factual, conceitual, procedimental e metacognitivo) classifica \emph{o que} o estudante
precisa aprender; a \textbf{dimensão dos processos cognitivos} classifica \emph{o que ele faz
mentalmente} com esse conhecimento. Para a Matemática do Ensino Médio, as categorias de conhecimento
mais relevantes são a conceitual e a procedimental.

A dimensão dos processos cognitivos organiza-se em seis categorias de complexidade crescente.
\textbf{Lembrar} corresponde à recuperação de informação armazenada; \textbf{entender}, a construir
significado (interpretar, exemplificar, comparar); \textbf{aplicar}, a executar um procedimento em uma
situação específica; \textbf{analisar}, a decompor o material e determinar como as partes se
relacionam; \textbf{avaliar}, a emitir julgamentos com base em critérios; e \textbf{criar}, a reunir
elementos para formar um todo novo e coerente. A hierarquia não é absoluta nem linear, mas indica uma
progressão geral que orienta o planejamento de avaliações equilibradas.

No contexto da Matemática do Ensino Médio, \citet{miranda2015taxonomia} mostrou que as questões
elaboradas por professores concentram-se maciçamente nos três primeiros níveis --- lembrar, entender e
aplicar ---, com presença muito reduzida de itens nos níveis superiores. Esse resultado converge com o
diagnóstico da BNCC, que aponta a necessidade de superar o modelo de avaliação centrado na reprodução
de procedimentos em direção ao desenvolvimento de competências de raciocínio e argumentação. Para o
sistema desta dissertação, a TBR fornece o vocabulário com que o professor declara o nível cognitivo
desejado; o detalhamento de como esse parâmetro opera no fluxo de execução é apresentado na
Seção~\ref{sub:sintese-tax}.

\subsection{Taxonomia SOLO}
\label{sub:solo}

A Taxonomia SOLO (\emph{Structure of Observed Learning Outcomes}) foi proposta por John Biggs e Kevin
Collis em 1982 \citep{biggs1982solo} para classificar a \emph{qualidade das respostas} produzidas por
estudantes. Enquanto a TBR parte da perspectiva do professor --- o que se pretende que o estudante
aprenda ---, a SOLO parte da perspectiva da resposta observada: dado o que o estudante produziu, qual é
a estrutura dessa produção? Essa diferença de orientação torna as duas taxonomias complementares, e não
concorrentes.

A SOLO organiza-se em cinco níveis. No nível \textbf{pré-estrutural}, o estudante não demonstra
compreensão relevante da tarefa. No \textbf{uniestrutural}, identifica um único aspecto relevante, mas
não o articula com outros. No \textbf{multiestrutural}, mobiliza vários aspectos, mas os trata de forma
independente, sem integração. No \textbf{relacional}, os diferentes aspectos são articulados em uma
estrutura coerente, de modo que o estudante demonstra compreensão integrada do problema. No
\textbf{abstrato ampliado}, vai além dos limites do problema dado, generalizando a estrutura para novos
contextos. A distinção decisiva é a que separa a passagem \emph{quantitativa} (do pré-estrutural ao
multiestrutural --- acumular elementos) da \emph{qualitativa} (do multiestrutural ao relacional ---
reorganizar os elementos em uma estrutura integrada).

Aplicada à elaboração de questões, a SOLO oferece um diagnóstico que o modelo de Bloom não captura com
a mesma clareza. Uma questão de função quadrática que pede apenas o cálculo das raízes pela fórmula de
Bhaskara admite uma resposta multiestrutural correta; uma questão que pede ao estudante que justifique
\emph{por que} o discriminante determina o número de raízes reais exige uma resposta relacional. Ambas
podem ser classificadas no nível \emph{aplicar} da TBR e ainda assim exigir estruturas de resposta
radicalmente diferentes --- e é essa diferença que a SOLO torna visível.

\subsection{Síntese operacional}
\label{sub:sintese-tax}

% CASA CANÔNICA do par "TBR na entrada, SOLO na avaliação". No rascunho
% esse mesmo argumento aparecia em 3.2.1, 3.2.2, 3.2.3 e 3.3.3; aqui ele
% é dito UMA vez e referenciado nos demais pontos.

A adoção conjunta das duas taxonomias responde a complementaridade funcional, não a exaustividade
teórica: a TBR responde a ``qual processo cognitivo a questão deve exigir?'' --- voltada ao \emph{design}
da tarefa ---; a SOLO responde a ``qual é a estrutura da resposta que a questão efetivamente
demanda?'' --- voltada à qualidade do produto. A escolha por essas duas, em vez de outras igualmente
reconhecidas, apoia-se em dois critérios: a compatibilidade com o vocabulário curricular brasileiro ---
a TBR já é referenciada implicitamente pela BNCC --- e o uso documentado de ambas em pesquisas
brasileiras com questões do Ensino Médio.

No sistema desta dissertação, essa complementaridade é operacionalizada em dois momentos do fluxo de
execução: a \textbf{TBR compõe a especificação de entrada} --- o professor seleciona um dos seis
processos cognitivos como parâmetro obrigatório, incorporado ao \emph{prompt} do Agente Gerador --- e a
\textbf{SOLO informa a avaliação} --- o Agente Crítico Didático verifica se a estrutura da resposta
esperada é compatível com o nível declarado. Uma questão especificada como \emph{analisar} deve exigir
resposta de nível relacional ou abstrato ampliado; se a resposta esperada for correta já no nível
multiestrutural, isso é sinalizado como inconsistência, e o orquestrador devolve a questão ao gerador
com instrução de revisão. Essa verificação cruzada --- TBR na entrada, SOLO na avaliação --- fecha o
ciclo de garantia de qualidade cognitiva do sistema.

\section{BNCC do Ensino Médio}
\label{sec:bncc}

Os referenciais anteriores fornecem os critérios pedagógicos e cognitivos da elaboração de questões.
Para torná-los operacionais, é preciso ancorá-los ao currículo concreto do Ensino Médio brasileiro ---
ancoragem fornecida pela Base Nacional Comum Curricular (BNCC), homologada em 2018
\citep{brasil2018bncc}, que organiza as habilidades de Matemática em um sistema de códigos que permite
ao professor especificar, de forma precisa, qual conteúdo e qual expectativa de aprendizagem estão em
jogo em cada questão gerada.

\subsection{Estrutura geral}
\label{sub:bncc-estrutura}

A BNCC do Ensino Médio organiza o currículo em quatro áreas do conhecimento, das quais interessa aqui a
de Matemática e suas Tecnologias. Cada área define \textbf{competências específicas} --- capacidades
amplas a serem desenvolvidas --- e as desdobra em \textbf{habilidades} concretas, identificadas por
códigos alfanuméricos. Além da formação básica comum, a BNCC prevê os \textbf{itinerários formativos};
no âmbito desta dissertação, o recorte recai sobre a formação básica comum da área de Matemática,
obrigatória para todos os estudantes, ficando os itinerários, por sua implementação heterogênea, fora
do escopo.

A estrutura tem uma consequência direta para o sistema: o código de habilidade --- da forma
\texttt{EM13MATxxx}, em que o primeiro algarismo após \texttt{MAT} identifica a competência específica e
os dois seguintes numeram a habilidade dentro dela --- funciona como identificador único e não ambíguo da expectativa de
aprendizagem associada a cada questão gerada, permitindo rastrear qual conteúdo ela aborda e como se
posiciona no currículo.

\subsection{Habilidades de Matemática}
\label{sub:bncc-habilidades}

Diferentemente do Ensino Fundamental, a área de Matemática e suas Tecnologias não distribui suas
habilidades por unidades temáticas: as quarenta e cinco habilidades da formação básica comum estão
organizadas pelas cinco competências específicas, e é o primeiro algarismo do código que indica a qual
delas cada habilidade pertence. A própria Base, contudo, registra que essa não é a única organização
possível e admite que, na elaboração de currículos, as habilidades sejam reagrupadas em unidades
temáticas, à maneira do Ensino Fundamental \citep{brasil2018bncc}. O documento ilustra esse
reagrupamento com uma unidade --- \textbf{Funções polinomiais de 1º e 2º graus} --- que reúne sete
habilidades provenientes de três competências distintas: \texttt{EM13MAT501} e \texttt{EM13MAT502}
(investigar padrões em tabelas e expressar algebricamente a generalização), \texttt{EM13MAT401} e
\texttt{EM13MAT402} (converter representações algébricas em geométricas no plano cartesiano),
\texttt{EM13MAT503} (investigar pontos de máximo e de mínimo), \texttt{EM13MAT507} (associar
progressões aritméticas a funções afins de domínio discreto) e \texttt{EM13MAT302} (resolver e
elaborar problemas modelados por essas funções).

Essa unidade exemplificada pela BNCC é adotada nesta dissertação como recorte curricular, na íntegra.
A ela acrescenta-se um segundo agrupamento, construído por analogia --- e não extraído do texto da Base
---, que reúne as habilidades correspondentes no campo do crescimento exponencial:
\texttt{EM13MAT303} (porcentagens e juros compostos), \texttt{EM13MAT304} (problemas com funções
exponenciais) e \texttt{EM13MAT508} (associar progressões geométricas a funções exponenciais de domínio
discreto). O recorte totaliza, portanto, dez habilidades, cobrindo funções afins, funções quadráticas,
funções exponenciais e as duas progressões.

A delimitação não é de conveniência: ela decorre do critério de verificabilidade que sustenta a
hipótese arquitetural desta dissertação (Seção~\ref{sec:hipotese}). Equações, funções e sequências
admitem formalização em um sistema de computação algébrica e, portanto, verificação determinística da
resposta, ao passo que habilidades de geometria plana e espacial, de contagem e probabilidade e de
estatística exigiriam núcleos de verificação de natureza inteiramente distinta, quando não recaem nos
limites de formalização discutidos na Seção~\ref{sub:cas-limites}. Pelo mesmo motivo ficam de fora, por
ora, as funções logarítmicas e trigonométricas, ainda que sejam vizinhas imediatas do recorte e
constituam a extensão mais natural do sistema. Como o recorte é curricular e não meramente
computacional, ele coincide com uma unidade que a própria Base reconhece como coerente --- o que
permite ao professor situar o que o sistema gera dentro de um bloco de ensino já constituído. A
tradução desse recorte em parâmetros de entrada do sistema é descrita na
Seção~\ref{sec:especificacao}.

\subsection{Articulação entre BNCC e taxonomias cognitivas}
\label{sub:bncc-taxonomias}

A BNCC e a TBR operam em níveis de abstração distintos, mas complementares: as habilidades da BNCC
especificam \emph{o que} o estudante deve aprender; a TBR, \emph{como} ele deve processar
cognitivamente esse conteúdo. No sistema, o código de habilidade e o nível cognitivo são dois
parâmetros independentes da especificação de entrada, e sua combinação é o que permite ao professor
solicitar, com precisão, uma questão sobre um conteúdo específico em um nível de complexidade
determinado. A articulação não é arbitrária: os verbos das próprias habilidades já carregam uma
expectativa de nível cognitivo --- \emph{interpretar} aponta para \emph{entender}; \emph{resolver e
elaborar problemas} combina \emph{aplicar} e \emph{analisar}; \emph{investigar} aponta para os níveis
superiores. Essa correspondência fornece a ancoragem com que o Agente Crítico Didático avalia a
coerência entre a habilidade declarada e o nível cognitivo da questão gerada (conforme operacionalizado
na Seção~\ref{sub:sintese-tax}). Essa dupla ancoragem --- curricular pela BNCC e cognitiva pela TBR --- é
o que confere ao sistema sua especificidade pedagógica.

\section{Critérios de qualidade para questões matemáticas}
\label{sec:criterios}

% CASA CANÔNICA dos seis critérios. É o ÚNICO lugar onde eles aparecem
% enumerados e formatados. As seções anteriores apenas os DERIVAM de
% cada referencial e remetem para cá.

As seções anteriores apresentaram quatro referenciais --- Pólya, a TSD, as taxonomias de Bloom e SOLO,
e as habilidades da BNCC ---, cada um contribuindo com um ângulo específico para a compreensão do que
torna uma questão matematicamente e pedagogicamente adequada. Esta seção consolida essas contribuições
em um conjunto coeso de seis critérios que operam simultaneamente como rubrica do Agente Crítico
Didático e como instrumento de avaliação pelo painel docente. Os critérios articulam-se em uma estrutura
que vai da dimensão mais técnica --- a correção matemática --- à mais pedagógica --- a relevância
formativa.

\begin{enumerate}
\item \textbf{Correção matemática.} Os dados são consistentes, o gabarito está correto e a resolução
não contém erros. É o critério mais fundamental e o único que admite avaliação binária; no sistema, é o
único avaliado de forma determinística, pelo Agente Verificador Simbólico, antes de a questão chegar ao
Agente Crítico. Questões cujo gabarito depende de dados omitidos, ou que admitem múltiplas
interpretações incompatíveis, são reprovadas neste critério, pois violam a exigência de Pólya de
solução determinada e verificável.

\item \textbf{Clareza do enunciado e ausência de ambiguidade.} Um enunciado claro permite ao estudante
identificar sem dúvida o que é dado, o que é pedido e quais são as condições --- a primeira etapa da
heurística de Pólya. À luz da TSD, a ambiguidade empobrece o \emph{milieu}: quando o estudante não
consegue determinar sobre qual problema deve raciocinar, a resistência produtiva do ambiente se dissolve
em confusão.

\item \textbf{Adequação ao nível cognitivo e ao público.} O processo cognitivo efetivamente demandado
deve ser coerente com o nível declarado e compatível com o repertório de um estudante do Ensino Médio.
Na dimensão cognitiva, o Agente Crítico verifica, com base na SOLO, se a estrutura de resposta exigida é
compatível com o nível Bloom declarado (Seção~\ref{sub:sintese-tax}); na curricular, verifica se os
conceitos envolvidos estão previstos para a etapa.

\item \textbf{Alinhamento à habilidade BNCC declarada.} A questão deve abordar o objeto de conhecimento
correspondente à habilidade declarada (plano do conteúdo) e ter demanda cognitiva compatível com a
expectativa expressa na habilidade (plano do verbo, conforme a Seção~\ref{sub:bncc-taxonomias}).

\item \textbf{Plausibilidade dos distratores} (em questões de múltipla escolha). Um distrator bem
elaborado é plausível: corresponde a um erro sistemático real --- confundir o discriminante com a raiz,
inverter numerador e denominador, esquecer um caso. Um distrator implausível não cumpre função
diagnóstica e reduz artificialmente a dificuldade.

\item \textbf{Originalidade e relevância pedagógica.} A questão não reproduz mecanicamente enunciados de
livros didáticos, e seu contexto é significativo para o estudante. À luz da TSD, a relevância está
ligada à qualidade do \emph{milieu}. É o critério mais subjetivo dos seis, avaliado em escala descritiva
de três pontos (inadequado, aceitável, rico), e não por veredicto binário.
\end{enumerate}

A Tabela~\ref{tab:criterios} sintetiza os seis critérios, seus referenciais de origem, o agente
avaliador e o tipo de resposta produzida.

\begin{table}[htbp]
\centering
\caption{Critérios de qualidade, referenciais e mecanismos de avaliação}
\label{tab:criterios}
\begin{tabular}{@{}p{3.2cm}p{2.4cm}p{2.8cm}p{3.4cm}@{}}
\toprule
\textbf{Critério} & \textbf{Referencial} & \textbf{Agente avaliador} & \textbf{Tipo de resposta} \\
\midrule
Correção matemática & Pólya; SymPy & Verificador Simbólico & Binário (aprovado / rejeitado) \\
Clareza do enunciado & Pólya; TSD & Crítico Didático & Escala + comentário \\
Adequação ao nível & Bloom; SOLO & Crítico Didático & Escala + comentário \\
Alinhamento BNCC & BNCC; Bloom & Crítico Didático & Binário + justificativa \\
Plausibilidade dos distratores & Dante; TSD & Crítico Didático & Escala por distrator \\
Originalidade e relevância & TSD; Pólya & Crítico Didático & Escala descritiva (3 pontos) \\
\bottomrule
\end{tabular}

\smallskip
{\footnotesize Elaborado pela autora.}
\end{table}

Os seis critérios formam a rubrica operacional do sistema e serão retomados no Capítulo~\ref{cap:avaliacao}
como instrumento de avaliação pelo painel docente: os professores participantes receberão uma versão
desta rubrica adaptada para avaliação humana, permitindo a comparação direta entre o veredicto do
Agente Crítico Didático e o julgamento dos especialistas. Essa triangulação --- entre verificação
simbólica determinística, avaliação automatizada por rubrica e avaliação humana qualificada --- é o
mecanismo de validação da hipótese arquitetural descrita na Seção~\ref{sec:hipotese}.
